La lumière jaune du sodium est formée de deux raies très proches, de longueurs d'onde \qty{589.0}{nm} et \qty{589.6}{nm}. On veut les séparer avec un réseau de \num{\v{Nmm}} traits par millimètre, dont un faisceau éclaire une largeur de \qty{\v{W}}{mm}. On observe le spectre d'ordre \v{m}.

\begin{enumerate}[label=\alph*)]
    \item Quel pouvoir de résolution faut-il pour séparer ces deux raies ?
    \item Combien de traits doivent être éclairés, au minimum, pour les séparer à l'ordre \v{m} ?
    \item Combien de traits sont éclairés ici ? Les deux raies sont-elles séparées ?
    \item Quel est le plus petit écart de longueur d'onde que ce montage peut résoudre près de \qty{589}{nm} ?
\end{enumerate}
